\section{제어기 설계}\label{sec:controllers}
PID 구현은 측정값에 대한 미분, 1차 미분 필터, 출력 포화, 역계산 안티와인드업을 사용한다. 저장된 적분 상태와 필터링된 미분 상태는 \SI{1}{s}마다 갱신하며 실제 시나리오의 출력 한계를 우선 적용한다.

LQR은
\begin{equation}
J=\sum_{k=0}^{\infty}(x_k^{\mathsf T}Qx_k+u_k^{\mathsf T}Ru_k) \label{eq:lqcost}
\end{equation}
를 최소화하고 이산 대수 Riccati 방정식에서 피드백 이득을 구한다. 추종에는 공칭 평형 앞먹임 $u=UA_0(r-T_{a0})-K(x-r)$를 사용한다. 따라서 일반 LQR은 미지의 일정 부하를 오프셋 없이 제거한다고 기대할 수 없다.

LQI는 다음과 같이 상태를 확장한다.
\begin{align}
z_{k+1}&=z_k+h(r_k-y_k),\nonumber\\
u_k&=UA_0(r-T_{a0})-K_x(x_k-r)-K_z z_k. \label{eq:lqi}
\end{align}
명령 계산에는 $z_{k+1}$ 저장 전의 $z_k$를 사용하여 합성 규약과 일치시켰다. 확장 시스템의 가제어성 행렬식은 $-hB_d^2\ne0$이다.

기존 PID 이득은 $(420,2,1200)$, LQR 가중치는 $(Q,R)=(1,2\times10^{-5})$, LQI 가중치는 $(Q_x,Q_z,R)=(1,8\times10^{-4},2\times10^{-5})$이다. 재튜닝 매개변수는 \cref{tab:tuned}과 같이 자동 생성한다.
\begin{table}[!htbp]\centering\caption{동일 예산 매개변수: PID $(K_p,K_i,K_d)$, LQR $(Q,R)$, LQI $(Q_x,Q_z,R)$.}\label{tab:tuned}\resizebox{\linewidth}{!}{\begin{tabular}{lrrr}\toprule
Controller & Parameter 1 & Parameter 2 & Parameter 3 \\
\midrule
PID & 596.5702 & 3.2242 & 2424.0000 \\
LQR & 1.821264e1 & 1.842874e-7 & -- \\
LQI & 1.400229e-1 & 1.315769e-3 & 6.155006e-6 \\
\bottomrule\end{tabular}
}\end{table}
